%%%%%%%%%%%%%%%%%%%%%%%%%%%%%%%%%%%%%%%%%%%%%%%%%%%%%%%%%%%%%%
\section{Predicados e Funções}


%%%%%%%%%%%%%%%%%%%%%%%%%%
\begin{frame}
\frametitle{Predicados e Funções}
\begin{minipage}{0.47\textwidth}
    \begin{itemize}
        \item Definições
        \item Contexto de uso
        \item Exemplos com a definições
        \item Estruturas de repetição
        \item Estruturas de decisão
    \end{itemize}
\end{minipage}
\begin{minipage}{0.5\textwidth}
\begin{figure}[ht!]
\begin{center}
\includegraphics[width=1.2\textwidth, height=0.40\textheight]{figures/logo_picat_alex.jpg}
\end{center}
\end{figure}
\end{minipage}
\end{frame}
%%%%%%%%%%%%%%%%%%%%%%%%%%


\begin{frame}[fragile,allowframebreaks=0.8,c]
	\frametitle{Predicados e Funções}
    
\begin{itemize}

\item Os \textbf{predicados} sempre assumem valores valores lógicos \texttt{true} ($1$) ou \texttt{false} ($0$).

\item Os \textbf{predicados} em seus argumentos, podem passar $n$-termos e receber
outros termos.

\item Quanto as \textbf{funções}, estas funcionam seguindo as regras de funções matemáticas,
sempre retornando um único valor

\item Predicados e funções são definidos com regras de \textit{casamento de padrões}


\item Predicados são conhecidos como \textbf{regras lógicas}, há dois tipos de regras:
\framebreak        
        \begin{itemize}
            \item Regras \textbf{sem} {\em backtracking} (\textit{non-backtrackable}):
        	
            \begin{tabbing}
              aa \= aaa \= aaa \= aaa \= aaa \= aaa \= aaa \kill
             \> \> \verb+ Cabeça , Condicional + \textbf{\textcolor{red}{=>}} \verb+ Corpo .+
            \end{tabbing}
            

            \item Regras \textbf{com} {\em backtracking}:
            
            \begin{tabbing}
              aa \= aaa \= aaa \= aaa \= aaa \= aaa \= aaa \kill
             \> \> \verb+ Cabeça , Condicional + \textbf{\textcolor{red}{?=>}} \verb+ Corpo .+
            \end{tabbing}
        
        \end{itemize}
\framebreak        
       
  \item A identificação da sintaxe é dada por:
  
  \begin{itemize}
  
      \item \emph{Cabeça}: indica um padrão de regra a ser casada.
      
      Forma geral:
      \begin{tabbing}
          aa \= aaa \= aaa \= aaa \kill
          \> \> $regra(termo_1, \ldots,termo_n)$
      \end{tabbing}
      
      Onde:
      
      \begin{itemize}
          \item \emph{regra} é um átomo que define o nome da regra.
          \item \emph{n} é a aridade da regra (\emph{i.e.} o total de argumentos)
          \item Cada \emph{$termo_i$} é um argumento da regra.
      \end{itemize}
      
      \item \emph{Cond}: é uma ou várias condições sobre a execução desta regra.
          
      \item \emph{Corpo}: define as ações da regra 
      
  \end{itemize}
  
  \item Todas as regras são finalizadas por um ponto final (.), seguido por um espaço 
  em branco ou nova linha.
  
  \item Ao longo dos exemplos, detalhes a mais sobre esta construção de \textbf{predicados} e \textbf{funções}
      
\end{itemize}

\end{frame}
%%%%%%%%%%%%%%%%%%%%%%%%%%%%%%%%%%%%%%%%%%%%%%%%%%%%%%%%%%%%%%%%%%%%%


\begin{frame}[fragile]
	\frametitle{Regras com e sem Backtracking -- Exemplo}


\begin{footnotesize}
\begin{verbatim}
main => regra_01(7)  ,
        regra_01(-4) ,
        regra_01(44) .
        
regra_01(0)          ?=> printf("\n  CHEGOU A 0 !!!\n").
regra_01(N) , N < 0  ?=> printf("\n  EH UM NEGATIVO !!!\n").
regra_01(N) , N > 10 ?=> printf("\n  EH MAIOR QUE 10 !!!\n").
regra_01(N) , N <= 10  =>  
          printf("\t :%d ", N),
          regra_01(N-1).

%% =<  EH IGUAL a >= :: sobrecarga
%%% $ picat regras_com_sem_backtraking.pi

\end{verbatim}

\end{footnotesize}
\textcolor{red}{Há uma recursão aqui. Algumas seções a frente.}

\end{frame}


%%%%%%%%%%%%%%%%%%%%%%%%%%%%%%%%%%%%%%%%%%%%%%%%%%%%%%%%%%%%%%%%%%%%%

\subsection{Casamento de Padrões}

\begin{frame}[fragile]
\frametitle{Casamento de Padrões}

\begin{itemize}

\item O algoritmo de {\em casamento de padrões} para regras é análogo ao algoritmo de    unificação para variáveis. 

\item O objetivo é encontrar dois padrões que possam ser unificados para se inferir alguma ação.

\pause
\item Muitos exemplos no curso

\end{itemize}
\end{frame}


%%%%%%%%%%%%%%%%%%%%%%%%%%%%%%%%%%%%%%%%%%%%%%%%%%%%%%%%%%%%%%%%%%%%%

\begin{frame}[fragile]
\frametitle{Casamento de Padrões -- Procedimento}


Procedimento de {\em casamento de padrões}:

\begin{itemize}
	\item Dado um padrão $p_1(t_1, \ldots,t_m)$, este \emph{casa} 
	com um padrão semelhante $p_2(u_1, \ldots,u_n)$ se:
    
    \begin{itemize}
    	\item $p_1$ e $p_2$ forem átomos equivalentes;
    
        \item O número de termos (chamado de aridade) em ($t_1, \ldots,t_m$)
        e ($u_1, \ldots,u_n$) for equivalente.
    
    	\item Os termos ($t_1, \ldots,t_m$) e ($u_1, \ldots,u_n$) são equivalentes, ou tornaram-se  equivalentes pela unificação de variáveis	em qualquer um dos dois termos;
    \end{itemize}
    
    \item Caso essas condições forem satisfeitas, o padrão $p_1(t_1,\ldots,t_m)$  casa com o padrão $p_2(u_1, \ldots,u_n)$.
    
\end{itemize}        

\end{frame}

%%%%%%%%%%%%%%%%%%%%%%%%%%%%%%%%%%%%%%%%%%%%%%%%%%%%%%%%%%%%%%%%%%%%%

\begin{frame}[fragile]
    
\begin{block}{Exemplos de \textit{casamento}}

\begin{enumerate}
	\item A regra $fatorial(Termo, Resultado)$ pode casar com:\\
    $fatorial(1,1)$, $fatorial(5,120)$, $fatorial(abc,25)$, $fatorial(X,Y)$, \emph{etc}.
    
    \pause
     \item A regra $fatorial(Termo, Resultado), Termo \geq 0$ pode casar com:\\
    $fatorial(1,1)$, $fatorial(5,120)$, $fatorial(X,Y)$, $fatorial(Z,Z)$, \emph{etc}.
    
    \pause
    \item A regra $pai(X,Y)$ pode casar com:\\
    $pai(rogerio, miguel)$, $pai(rogerio, henrique)$, $pai(salomao, X)$, $pai(12,24)$, \emph{etc}.
    
    \pause
    \item A regra $pai(salomao, X)$ pode casar com:\\
    $pai(salomao, rogerio)$, $pai(salomao, fabio)$.
    
    \pause
    \item A regra $pai(salomao, fabio)$ pode casar com:\\
    $pai(X, fabio)$, $pai(salomao, X)$, $pai(X,Y)$
\end{enumerate}

\end{block}
    
\end{frame}

%%%%%%%%%%%%%%%%%%%%%%%%%%%%%%%%%%%%%%%%%%%%%%%%%%%%%%%%%%%%%%%%%%%%%

\begin{frame}[fragile]

	\frametitle{Metas ou Provas -- (\textit{goals})}
	
	\begin{itemize}
	    
	 \item Na matemática ao se deduzir um valor de um teorema, tem-se uma \textit{prova}. 
	  Assim, o termo  \textit{goal} eventualmente é chamado de \textit{prova do programa}
	 
	 \pause
	    \item Metas ou Provas (do inglês: \textit{goal}) são estados que definem o final da execução.
	    
  	 \pause
	    \item Uma meta pode ser, entre outros, um valor lógico, uma chamada de outra regra, 
	    uma exceção ou uma operação lógica.
	\end{itemize}

\end{frame}
    
  %  \begin{enumerate}
  %      \item \texttt{true, yes} $\Rightarrow$ Valor lógico para verdade.
%   %      \item \texttt{false, no} $\Rightarrow$ Valor lógico para falsidade.
%         \item \texttt{$p(t_1,\ldots,t_n)$} $\Rightarrow$ Chamada de uma regra \emph{p}.
%         \item \texttt{(P,Q), (P;Q), (P\&\&Q), (P||Q), not P} $\Rightarrow$ Operação lógica 
%     %    sobre uma ou mais metas P e Q.
   % \end{enumerate}
%%%%%%%%%%%%%%%%%%%%%%%%%%%%%%%%%%%%%%%%%%%%%%%%%%%%%%%%%%%%%%%%%%%%%
\subsection{Funções}

\begin{frame}[fragile]
	\frametitle{Funções -- I}
    
    \begin{itemize}
        \item A forma geral de uma função é:
        
        \begin{tabbing}
          aa \= aaa \= aaa \= aaa \= aaa \= aaa \= aaa \kill
          \> \> $Cabe$ç$a = X \ $\verb+=>+$\ Corpo$. 
        \end{tabbing}

       \pause        
        \item Caso tenhamos uma condição \emph{Cond}::
        
        \begin{tabbing}
          aa \= aaa \= aaa \= aaa \= aaa \= aaa \= aaa \kill
          \> \> $Cabe$ç$a = X , Cond \ $\verb+=>+$\ Corpo$. 
        \end{tabbing}
        
        \item \textcolor{red}{Funções \textbf{não} admitem \textbf{\textit{backtracking}}}.
    \end{itemize}
\end{frame}
    
\begin{frame}[fragile]
	\frametitle{Funções -- II}

    
    \begin{itemize}
        \item Funções são tipos especiais de regras que sempre sucedem com \emph{uma} 
        resposta.
        
        \pause
        \item Funções em Picat tem como intuito serem sintaticamente semelhantes a funções matemáticas (vide \emph{Haskell}).
        
        \pause
        \item Em uma função a \emph{Cabeça} é uma equação do tipo $f(t_1,\ldots,t_n)=X$, onde
        \emph{f} é um átomo que é o nome da função, \emph{n} é a aridade da função, e cada
        termo $t_i$ é um argumento da função.
        
        \item \emph{X} é uma expressão que é o retorno da função.
        
    \end{itemize}
    
\end{frame}
    
\begin{frame}[fragile]
	\frametitle{Funções  -- III}
    
    \begin{itemize}
    
        \item Funções também podem ser denotadas como fatos, onde podem servir como 
        \textbf{\textit{aterramento}} para regras recursivas.
        %, ou até mesmo como versões simplificadas de uma      regra. 
        
        
        \item Esta são denotadas como: $f(t_1,\ldots,t_n) = $ \emph{Expressão}, onde \textit{Expressão}
        pode ser um valor ou uma série de ações.
        
    \end{itemize}
    
\end{frame}

%%%%%%%%%%%%%%%%%%%%%%%%%%%%%%%%%%%%%%%%%%%%%%%%%%%%%%%%%%%%%%%%%%%%%
\begin{frame}[fragile]
\frametitle{Regras, Metas e Funções -- Exemplo}

\begin{footnotesize}
\begin{verbatim}
main =>
       X = 3, 
       Y = 4,
       um_predicado(X,Y,Z),
       R = uma_funcao(X,Y),
       printf("\n Z: %d \t R: %d", Z, R),
       println("\n FIM"). 
       
um_predicado(X,Y,Z) => Z = X + Y.

uma_funcao(X,Y) = R => R = X + Y.
\end{verbatim}

\end{footnotesize}
\end{frame}


\begin{frame} [fragile]

\frametitle{Regras, Metas e Funções -- Exemplo}
\begin{footnotesize}
\begin{verbatim}
Picat> cl('predicados_funcoes').
Compiling:: predicados_funcoes.pi
predicados_funcoes.pi compiled in 0 milliseconds
loading...

yes

Picat> um_predicado(3,4,Z), write(Z).
7Z = 7
yes

Picat> uma_funcao(3,4) = R, write(R).
7R = 7
yes

Picat>
\end{verbatim}
\end{footnotesize}
\end{frame}

%%%%%%%%%%%%%%%%%%%%%%%%%%%%%%%%%%%%%%%%%%%%%%%%%%%%%%%%%%%%%%%%%%%%%

\subsection{Relembrando as Regras}
\begin{frame}[c, fragile, allowframebreaks=0.75]
	\frametitle{Relembrando as Regras}
    
    \begin{itemize}
        \item Forma geral de um predicado do tipo Regra:
    
        \begin{tabbing}
          aa \= aaa \= aaa \= aaa \= aaa \= aaa \= aaa \kill
          \> \> \verb+ Cabeça , Condicional  =>   Corpo .+ 
        \end{tabbing}
        
        \item  Forma geral de um predicado com \textit{backtracking}:
        
        \begin{tabbing}
          aa \= aaa \= aaa \= aaa \= aaa \= aaa \= aaa \kill
            \> \> \verb+ Cabeça , Condicional + \textbf{\textcolor{red}{?=>}} \verb+ Corpo .+
        \end{tabbing}
    \end{itemize}
    
    \textcolor{red}{Em Prolog, esta condicional (Cond) entra no corpo da regra. Picat é flexível!}

    \framebreak
    
    \begin{itemize}
        \framebreak
        
        \item Dentro de uma regra, \emph{Cond} só pode ser avaliado uma vez, acessando somente termos dentro do escopo do predicado.
        
        \item Sempre estar atento que: regras são \textbf{sempre avaliados com valores lógicos} ({\em true} ou {\em false})
        
        \item Por outro lado, as variáveis  como argumento ou  instanciadas
        dentro dele, podem ser utilizadas dentro do escopo da regra, 
        ou no escopo onde  esta regra foi chamada.
    \end{itemize}
    
\end{frame}

%%%%%%%%%%%%%%%%%%%%%%%%%%%%%%%%%%%%%%%%%%%%%%%%%%%%%%%%%%%%%%%%%%%%%

\subsection{Regras do Tipo  Fatos}
\begin{frame}[c,allowframebreaks=0.6,fragile]
\frametitle{Regras do Tipo  Fatos}

    \begin{itemize}
        \item As regras que não tem condicionais e nem corpo, estes  são conhecidos
        como:  \textit{fatos} ou \textit{regras sem-corpo}
        
       \item Estes \textit{fatos} são regras \textit{sempre verdadeiras}
        
        \item Formato dos fatos são do tipo:
        
        \begin{tabbing}
          aa \= aaa \= aaa \= aaa \= aaa \= aaa \= aaa \kill
          \> \> $p(t_1,\ldots,t_n)$. 
        \end{tabbing}
        
        \item Os argumentos de um \emph{fato} \textbf{não} podem conter \textbf{variáveis}.
        
        \framebreak
        
        \item A declaração de um fato é precedida por uma declaração \textbf{\emph{index}},
        algo como:
        
        \begin{tabbing}
        aa \= aaa \= aaa \= aaa \= aaa \= aaa \= aaa \kill
            \> \texttt{index ($M_{11},M_{12},\ldots,M_{1n}$) $\ldots$ ($M_{m1},M_{m2},\ldots,M_{mn}$)}\index{\texttt{index}} 
        \end{tabbing}
        
        \item Onde um $M_{ij}$ com o simbolo `$+$',  significa que este termo já foi 
        indexado.
        
        \item Quanto o `$-$' significa que este termo deve ser indexado
        
        \item Ou seja, quando ocorre um simbolo `$+$' em um grupo do \textit{index}, é avaliado
        pelo compilador como um valor constante a ser casada
        
        %, que não irá gerar uma nova regra durante
        %a execução do programa.
        
        \item Quanto ao `$-$', este é avaliado pelo compilador com uma variável que deverá ser 
        instanciada à um valor. Ou seja, quando se deseja unificar um valor a esta variável
        
        \item Dica: o parâmetro `$-$' no \textit{index} é quase como regra geral
                
        \item \textbf{Não} pode haver um \textbf{predicado} e um \textbf{predicado 
        fato} com \textbf{mesmo nome}.
        
    \end{itemize}
\end{frame}



\begin{frame} [fragile]
\frametitle{Exemplo -- Função e Regras}

\begin{footnotesize}
\begin{verbatim}
index (+,+,+) (+,+, -) (-,+,-) (-,-,+) (-,-, +) 
      (+,-,+) (+,+,-)  (-,-,-)
%(-,-,-) %% NENHUM argumento instanciao -- UTIL
%(+,+,+) %% TODOS ARGUMENTOS DEVEM ESTAR INSTANCIADOS
%(+,+, -) (-,+,-) (-,-,+) (-,-, +) (+,-,+) (+,+,-) (-,-,-)

and2(true,true,true).
and2(true,false,false).
and2(false,true,false).
and2(false,false,false).

main ?=>
       and2(X,Y,Z), % and eh reservado
       printf("\n X: %w \t Y: %w \t Z: %w", X, Y, Z),
       fail.
main =>       
       println("\n FIM"). 
\end{verbatim}
\end{footnotesize}

\textcolor{red}{Este exemplo é muito interessante. Execute ele na console do interpretador
excluindo alguns dos parâmetros do \textit{index}}

\end{frame}


%%%%%%%%%%%%%%%%%%%%%%%%%%%%%%%%%%%%%%%%%%%%%%%%%%%%%%%%%%%%%%%%%%%%%


%%%%%%%%%%%%%%%%%%%%%%%%%%%%%%%%%%%%%%%%%%%%%%%%%%%%%%%%%%%%%%%%%%%%%

\subsection{Exemplos}

\begin{frame}[fragile]

\frametitle{Exemplos}

\begin{block}{Exemplo de Predicado ou regra}
   
\begin{lstlisting}[frame=single]
contas_P0(X1, X2, X3, Z) ?=>
    number(X1),
    number(X2),
    number(X3),
    X1 < X2,
    X2 < X3,
    Z  =  (X2 + X3).
    
contas_P0(X1, _, _, Z) =>
   Z = X1.    
\end{lstlisting}
    \end{block}

        
\end{frame}

\begin{frame}[fragile]

\begin{block}{Exemplo de Funções}
     
\begin{lstlisting}[frame=single]
contas_F0(X1, X2, X3) = Z, (number(X1), 
                           number(X2), 
                           number(X3)) =>
    if (X1 < X2 && X2 < X3) then
    Z  =  (X2 + X3)
    else
    Z  =  X1
    end.
\end{lstlisting}
        
 \end{block}

\textcolor{red}{\textit{Aperitivo} à próxima seção: condicionais e laços!}
    
\end{frame}    
\begin{frame}[fragile]
   
\begin{block}{Mais Exemplos (Fatos e Regras)}
    
\begin{lstlisting}[frame=single]
index(-,-) (+,-) (-,+)
pai(salomao, rogerio).
pai(salomao, fabio).
pai(rogerio, miguel).
pai(rogerio, henrique).

avo(X,Y) ?=> pai(X,Z), pai(Z,Y).
irmao(X,Y) ?=> pai(Z,X), pai(Z,Y).
tio(X,Y) ?=> pai(Z,Y), irmao(X,Z).
\end{lstlisting}
    
 \end{block}
    
\end{frame}    
%%%%%%%%%%%%%%%%%%%%%%%%%%%%%%%%%%%%%%
\begin{frame}[fragile]
    
\begin{block}{Exemplos de Funções -- Equivalentes}
    
\begin{lstlisting}[frame=single]
eleva_cubo(1) = 1.
eleva_cubo(X) = X**3.
eleva_cubo(X) = X*X*X.
eleva_cubo(X) = X1 => X1 = X**3.
eleva_cubo(X) = X1 => X1 = X*X*X.
\end{lstlisting}
        
\end{block}
    
\end{frame}

%%%%%%%%%%%%%%%%%%%%%%%%%%%%%%%%%%%%%%%%%%%%%%%%%%%%%%%%%%%%%%%%%%%%%
%%% ATE AQUI
\section{Comandos Condicionais, Laços e Repetições}

\begin{frame}[fragile]
\frametitle{Comandos Condicionais, Laços e Repetições}
    
  
  \begin{itemize}
      
      \item Ao contrário do Prolog, Picat apresenta
      conceitos e comandos da programação imperativa
      
      \pause
      \item Esta maneira  ameniza os obstáculos
      em se aprender uma linguagem com o paradigma lógico,
      tendo outros elementos conhecidos
  
          \pause
      \item Assim, Picat apresenta estruturas clássicas
      como: 
      \begin{itemize}
        \item  \texttt{if--then--end}, \texttt{if--then--else--end}, \texttt{if--then--elseif--then--....end}
        \item  \texttt{foreach}
        \item  \texttt{while}
        \item  \texttt{do--while}
        \item Bem como a atribuição, `\texttt{:=}', já discutida
      \end{itemize}
   \end{itemize}

       
\end{frame}    
     
    %%%%%%%%%%%%%%%%
    
\begin{frame}[fragile]
\frametitle{Comandos Condicionais}
    
    \begin{itemize}
        
        \item Picat implementa uma  estrutura condicional explícita (na programação em lógica, voce faz isto implicitamente)
        
        \pause
        \item Sua notação é:\\
        
        \begin{tabbing}
            aa \= aaa \= aaa \= aaa \= aaa \= aaa \= aaa \kill
            \> \texttt{if (Exp) then}\\
            \> \> \texttt{Ações} \\
            \> \texttt{else}\\
            \> \> \texttt{Ações}\\
            \> \vdots\\
            \> \texttt{end} 
        \end{tabbing}
        
        \pause
        \item Onde \emph{Exp} é uma expressão lógica  avaliada como verdadeira ou  falsa.
     
     \pause   
        \item A última ação antes de um \emph{else} ou \emph{end} não deve ter  vírgula nem ponto e vírgula ao final da linha.
  
     \pause   
        \item Tem-se ainda o \texttt{elseif} que pode estar embutido no comando \texttt{if--then--else--end}
        
    \end{itemize}
\end{frame}    
    
    
    
    
    
    
    
    
\begin{frame}[fragile]

\begin{block}{Exemplo: \texttt{if--then--else--end} }
     
\begin{lstlisting}[frame=single]
    ,
if (X <= 100) then
        println("X e menor que 100")
    elseif (X <= 1000 && X >= 500) then
        println("X estah entre 500 e 1000")
    else
        println("X estah abaixo de  500")
    end
    ,
\end{lstlisting}
        
\end{block}
\end{frame}    

    
    
%%%%%%%%%%%%%%%%%%%%%%%%%%%%%%%%%%%%%%%%%%%%%%%%%%%%%%%%%%%%%%%    
    
    
    
\begin{frame}[fragile]
\frametitle{Estruturas de Repetições}
       

\begin{itemize}
    
    \item Picat também implementa 3 estruturas de repetição, são elas:
    \texttt{foreach}, \texttt{while}, e \texttt{do-while}.
    
    \pause
    \item O laço do  \texttt{foreach}  itera sobre termos simples e compostos.
    
        \pause
    \item O  \texttt{while} repete um conjunto de ações enquanto uma 
     condição for verdadeira.
     
        \pause
    \item A condição pode ser simples ou combinada
     
     
        \pause
    \item O laço \texttt{do-while} é análogo ao  \texttt{while}, porém ele 
    sempre executa pelo menos uma vez.
\end{itemize}
\end{frame}    


%%%%%%%%%%%%%%%%%%%%%%%%%%%%%%%%%%%%%%%%%%%%%%%%%%%%%%%%%%%%%%%    
\begin{frame}[fragile]
\frametitle{Estruturas de Repetições: \texttt{foreach}}

    \begin{itemize}
        \item Um laço \texttt{foreach} tem a seguinte forma:
        
        \begin{tabbing}
            aa \= aaa \= aaa \= aaa \= aaa \= aaa \= aaa \kill
            \> \texttt{foreach ($E_1$ in $D_1$, $Cond_1$, $\ldots$, $E_n$ in $D_n$, $Cond_n$)}  \\
            \> \> $Metas$ \\
            \>  \texttt{end} 
        \end{tabbing}
    
   \pause    
   Esta notação é dada por:
        
    \item  $E_i$ é um \emph{padrão de iteração} ou \emph{iterador}. 
        
    \item  $D_i$ é  uma expressão de \emph{valor composto}. Exemplo: uma lista de valores
        
     \item  $Cond_i$ é uma condição opcional sobre os iteradores $E_1$ até $E_i$.
        
     \item Laços do \texttt{foreach} podem conter múltiplos iteradores. 
     Caso isso ocorra, o compilador interpreta isso como diversos laços aninhados.
        
        % Maiores detalhes, ver Manual do Usuário. % copiar para outros slides
        
        % Exemplo 1 iterador e multiplos iteradores
        
    \end{itemize}
\end{frame}    


\begin{frame}[fragile]

\begin{block}{Exemplo: \texttt{foreach} }
     
\begin{lstlisting}[frame=single]
    
laco_01 =>  
    L = [17, 3, 41, 25, 8, 1, 6, 40],  
    foreach (E in L)  
        println(E)  
    end.  
\end{lstlisting}
        
\end{block}
\end{frame}    


%%%%%%%%%%%%%%%%%%%%%%%%%%%%%%%%%%%%%%%%%%%%%%%%%%%%%%%%%%%%%%%    
\begin{frame}[fragile]
\frametitle{Estruturas de Repetições: \texttt{while}}


\begin{itemize}
        \item O laço do \texttt{while} tem a seguinte forma:
        
        \begin{tabbing}
            aa \= aaa \= aaa \= aaa \= aaa \= aaa \= aaa \kill
            \> \texttt{while ($Cond$)} \\
            \> \> $Metas$  \\
            \>  \texttt{end}
        \end{tabbing} 
        
        \item Enquanto a expressão lógica \emph{Cond} for verdadeira, o conjunto de \texttt{Metas} é executado.
        
\end{itemize}
\end{frame}    


\begin{frame}[fragile]

\begin{block}{Exemplo: \texttt{while} }
     
\begin{lstlisting}[frame=single]
  
laco_02 =>  
    I = 1,  
    while (I <= 9)  
        println(I),  
        I := I + 2  
    end.  
\end{lstlisting}
        
\end{block}
\end{frame}    

%%%%%%%%%%%%%%%%%%%%%%%%%%%%%%%%%%%%%%%%%%%%%%%%%%%%%%%%%%%%%%%    
\begin{frame}[fragile]
\frametitle{Estruturas de Repetições: \texttt{do-while}}

    \begin{itemize}
        
        \item O laço \texttt{do-while} tem a seguinte forma:
        
        \begin{tabbing}
            aa \= aaa \= aaa \= aaa \= aaa \= aaa \= aaa \kill
            \> \texttt{do} \\
            \> \> $Metas$  \\
            \>  \texttt{while ($Cond$)}
        \end{tabbing} 
        
        \item Ao contrário do  \texttt{while} o iterador \texttt{do-while}
        vai executar \texttt{Metas} pelo menos uma vez antes de avaliar \texttt{Cond}.
        
    \end{itemize}
    
\end{frame}


\begin{frame}[fragile]

\begin{block}{Exemplo: \texttt{do-while} }

\begin{lstlisting}[frame=single]
     
laco_03 =>  
     J = 6,  
    do  
        println(J),  
        J := J + 1  
    while (J <= 5).  
\end{lstlisting}
        
\end{block}
\end{frame}    


%%%%%%%%%%%%%%%%%%%%%%%%%%%%%%%%%%%%%%%%%%%%%%%%%%%%%%%%%%%%%%%%%%%%%

\subsection{Funções e Predicados Especiais}


\begin{frame}[fragile]
%[c,allowframebreaks]

\frametitle{Funções e Predicados Especiais}
    % Compreensão de lista
    % Entrada e Saida de Dados
\begin{itemize}
    
    \item Há algumas funções e predicados especiais em Picat 
    que necessitam de algum cuidado.
    
    \pause
    \item São elas: \texttt{compreensão de listas/vetores}, 
    \texttt{entrada de dados} e \texttt{saída de dados}.
    
    \item Na verdade, já fizemos uso delas, porém sem a ênfase de que
    são funções ora predicados.
    
\end{itemize}
\end{frame}      


%%%%%%%%%%%%%%%%%%%%%%%%%%%%%%%%%%%%%%%%%%%%%%%%
\begin{frame}[c,allowframebreaks]

\frametitle{Compreensão de Listas e Vetores}    
    
\begin{itemize}    
    \item A função de \underline{\textit{compreensão de listas e vetores}} é uma função especial que  permite a fácil criação de listas ou vetores, opcionalmente seguindo uma regra de 
    criação.
    
    \item Sua notação é:
    
    \begin{tabbing}
        aa \= aaa \= aaa \= aaa \= aaa \= aaa \= aaa \kill
        \> \> \texttt{[$T$ : $E_1$ \texttt{in} $D_1$, $Cond_1$, $\ldots$, $E_n$ in $D_n$, $Cond_n$]} 
    \end{tabbing}
    
    \item Onde, $T$ é uma expressão  adicionada a lista, cada $E_i$ é um 
    iterador, cada $D_i$ é um termo composto ou expressão que gera um termo composto, e cada $Cond_i$ é uma condição sobre cada iterador de $E_1$ até $E_i$.
    
    \item Há uma seção dedicada a listas. Voltaremos ao assunto.
    
    \framebreak

    \item Esta função pode gerar um vetor também, a notação é um pouco diferente:
    
    \begin{tabbing}
        aa \= aaa \= aaa \= aaa \= aaa \= aaa \= aaa \kill
        \> \> \texttt{\{$T$ : $E_1$ \texttt{in} $D_1$, $Cond_1$, $\ldots$, $E_n$ in $D_n$, $Cond_n$\}} 
    \end{tabbing}
    
    \item Neste caso, os delimitadores são $\{$ e $\}$ de um vetor

\end{itemize}

\end{frame}      

% %%%%%%%%%%%%%%%%%%%%%%%%%%%%%%%%%%%%%%%%%    
 \begin{frame}[fragile]
 \frametitle{Compreensão de Listas e Vetores: Exemplo}
\begin{footnotesize}

\begin{verbatim}
main =>
 L = [(A, I) : A in [a, b], I in 1 .. 2],
 V = {(I, A) : A in [a, b], I in 1 .. 2},
 printf("\nL: %w  \nV: %w\n", L, V),
 imp_vetor(V).

imp_vetor (M) => 
Tam = M.length,  %% tamanho de M
  nl,
   foreach(I in 1  .. Tam )
     printf("V(%d):%w \t" , I, M[I] )
    end,
  nl.
%%%% $picat vetor_exemplo_01.pi
\end{verbatim}

\end{footnotesize}
 
\end{frame}


%%%%%%%%%%%%%%%%%%%%%%%%%%%%%%%%%%%%%%%%%%%%%%%%%%%%%%
\begin{frame}[c,allowframebreaks]

\frametitle{Leitura e Escrita}    

\begin{itemize}
        
\item Picat tem diversas variações  funções de leitura de valores,
        que serve tanto para ler de uma console \texttt{stdin},
        como de um arquivo qualquer.
        
\item Aos usuários de Prolog, aqui não precisamos do delimitador
final de `\textbf{.}'   ao final de uma leitura. 

\item Válido quando editamos no interpretador, o `\textbf{.}' final é opcional
        
\framebreak        
\item As mais importantes são:

\begin{itemize}
  
  \item \texttt{read\_int($FD$) = $Int$} $\Rightarrow$ Lê um \textit{Int} do 
  arquivo $FD$.
  
  \item \texttt{read\_real($FD$) = $Real$} $\Rightarrow$ Lê um \textit{Float} do 
  arquivo $FD$.
  
  \item \texttt{read\_char($FD$) = $Char$} $\Rightarrow$ Lê um \textit{Char} do 
  arquivo $FD$.
  
  \item \texttt{read\_line($FD$) = $String$} $\Rightarrow$ Lê uma \textit{Linha} 
  do arquivo $FD$.
  
\end{itemize}

\item Caso se deseja ler da console, padrão \texttt{stdin}, $FD$, o nome do descritor de arquivo,  pode ser omitido.

\framebreak

\item Os dois predicados mais importantes para saída de dados, são 
\texttt{write} e \texttt{print}.

\item Cada um destes predicados tem três variantes, são eles:

\begin{itemize}

\item \texttt{write($FD$, $T$)} $\Rightarrow$ Escreve um termo $T$ no arquivo 
$FD$.

\item \texttt{writeln($FD$, $T$)} $\Rightarrow$ Escreve um termo $T$ no arquivo 
$FD$, e pula uma linha ao final do termo.

\item \texttt{writef($FD$, $F$, $A\ldots$)} $\Rightarrow$ Este predicado é usado 
para escrita formatada para um arquivo $FD$, onde $F$ indica uma série de 
formatos para cada termo contido no argumento $A\ldots$. O número de argumentos 
não pode exceder 10.

\end{itemize}
        
\framebreak
        
\item Analogamente, para o predicado \texttt{print}, temos:

\begin{itemize}

  \item \texttt{print($FD$, $T$)} $\Rightarrow$ Escreve um termo $T$ no arquivo 
  $FD$.
  
  \item \texttt{println($FD$, $T$)} $\Rightarrow$ Escreve um termo $T$ no arquivo 
  $FD$, e pula uma linha ao final do termo.
  
  \item \texttt{printf($FD$, $F$, $A\ldots$)} $\Rightarrow$ Este predicado é usado 
  para escrita formatada para um arquivo $FD$, onde $F$ indica uma série de 
  formatos para cada termo contido no argumento $A\ldots$. O número de argumentos 
  não pode exceder 10.

\end{itemize}

\item Caso queira escrever para \texttt{stdout}, o nome do  $FD$, pode ser omitido.

\end{itemize}
    
\end{frame}

%%%%%%%%%%%%%%%%%%%%%%%%%%%%%%%%%%%%%%%%%%%%%%%%%%%%%%%%%%%%%%%%%%%%%

\begin{frame}
{Tabela de Formatação para Escrita}

Apenas os mais importantes, há outros como: hexadecimal, notação científica,
etc. Ver no apêndice do Guia do Usuário.    
\begin{table}[h]
    \begin{tabular}{l|l}
        \hline  \hline
        \textbf{Especificador} & \textbf{Saída} \\
        \hline 
        \hline 
        \texttt{\%\%} & Sinal de Porcentagem \\
        \texttt{\%c} & Caráctere \\
        \texttt{\%d \%i} & Número Inteiro Com Sinal \\
        \texttt{\%f} & Número Real \\
        \texttt{\%n} & Nova Linha \\
        \texttt{\%s} & \textit{String} \\
        \texttt{\%u} & Número Inteiro Sem Sinal \\
        \texttt{\%w} & Termo \textbf{qualquer} \\
        \hline   \hline
    \end{tabular}
\end{table}
\end{frame}

%%%%%%%%%%%%%%%%%%%%%%%%%%%%%%%%%%%%%%%%%%%%%%%%%%%%%%%%%%%%%%%%%%%%%

\begin{frame}

\frametitle{Comparação entre \texttt{write} e \texttt{print}}
     
\begin{table}[h]
    \centering
    \begin{tabular}{ l|l|l|l }
        \hline         \hline
        Dados $\Rightarrow$  & "abc" & [a,b,c] & 'a@b'\\
        \hline         \hline
        \texttt{write}   & [a,b,c] & [a,b,c] & 'a@b' \\
        \hline
        \texttt{writef}  & [a,b,c] (\%s) & abc (\%w) & 'a@b' (\%w) \\
        \hline
        \texttt{print}   & abc & abc & a@b \\
        \hline
        \texttt{printf}  & abc (\%s) & abc (\%w) & a@b (\%w) \\ 
        \hline \hline
    \end{tabular}
\end{table}

\end{frame}

%%%%%%%%%%%%%%%%%%%%%%%%%%%%%%%%%%%%%%%%%%%%%%%%%%%%%%%%%%%%%%%%%%%%%

\begin{frame}[fragile]

\frametitle{Exemplos}
    
    \begin{block}{Condicionais}
        \begin{lstlisting}[frame=single]
main =>
    X = read_int(),
    if(X <= 100)then
        println("X e menor que 100")
    else
        println("X nao e menor que 100")
    end.
.
        \end{lstlisting}
    
    \end{block}
\end{frame}

\begin{frame}[fragile]

\frametitle{Exemplos -- Repetições}
    
   \begin{lstlisting}[frame=single]
main =>
    X = read_int(),
    println(x=X),
    while(X != 0)
        X := X - 1,
        println(x=X)
    end
.
   \end{lstlisting}
        
   \begin{lstlisting}[frame=single]
main =>
    X = read_int(),
    Y = X..X*3,
    foreach(A in Y)
        println(A)
    end.
        \end{lstlisting}
    
\end{frame}

\begin{frame}[fragile]
\frametitle{Exemplos -- Compreensão de Listas}    
       
 \begin{lstlisting}[frame=single]
main => 
  Tamanho = read_int(),
  Y = [read_int() : 1..Tamanho],
  Z = {X : X in Y, X >= 0},
  forach(I in 1..Tamanho)
      if(Y[I] >= 0)then
          printf("Y[%d] e maior ou igual que 0\n",I)
      else
          printf("Y[%d] e menor que 0\n",I)
      end
  end
 printf("Os valores de Y que sao maiores que 0 sao:\n%w",Z)
.
\end{lstlisting} 
       
\textcolor{red}{Este exemplo reúne muitos conceitos desta seção.}       
   
\end{frame}

%%%%%%%%%%%%%%%%%%%%%%%%%%%%%%%%%%%%%%%%%%%%%%%%%%%%%%%%%%%%%%%%%%%%%

\begin{frame}[fragile]
\frametitle{Reflexões}

\begin{itemize}


  \item Esta seção  trata da sintaxe do Picat
  
    \pause
  \item Embora sua sintaxe não seja muito extensa, ela precisa 
  ser praticada
  
    \pause
  \item Como este conteúdo se assemelha as LPs clássicas, como exercício,
  voce está apto a fazer alguns algoritmos de outras linguagens.

  \pause
  \item Se antes o Prolog era complicado, com Picat,
  tudo ficou análogo a Python e a linguagem C
  
  \pause
  \item Em \url{https://github.com/claudiosa/CCS/tree/master/picat/input_output_exemplos}
  tem vários exemplos avançados de entradas e saídas 
  
  \pause
  \item Mãos à obra!
 
\end{itemize}

\end{frame}
